\section{Preliminaries} \label{sec:prelim}

\paragraph{General}
As usual, $\N, \R, \C $ stand for the natural, real, and complex numbers, respectively. %Given $x\in\C$, $x^*$ denotes its complex conjugate.
%
We adopt the Dirac notation for vectors representing quantum states, e.g., $\ket\psi, \ket\phi$, etc. In this paper, all the vectors of the form $\ket\psi$ are unit vectors. Given any pure state $\ket\psi$, we adopt the convention that its density operator is denoted by the Greek letter without the ``ket'', e.g. $\psi = \ket\psi\bra\psi$. The set of density operators in an arbitrary Hilbert space $\cH$ is denoted $\cD(\cH)$.
A symmetric state $\ket\psi \in (\C^d)^{\otimes n}$ is that invariant under any permutation $\pi\in \sym_n$ where $\sym_n$ is the symmetric group. The action of $\pi$ on $(\C^d)^{\otimes n}$ is  
\[
\pi:|\psi_1, \psi_2, \ldots, \psi_n\rangle
\mapsto |\psi_{\pi(1)}, \psi_{\pi(2)}, \ldots, \psi_{\pi(n)} \rangle.
\]
The \emph{symmetric subspace} is the subspace of $(\C^d)^{\otimes n}$ that is invariant under $\sym_n$, denoted by $\Symm{n}{d}$. Given any set $H\subseteq \cH$ for some Hilbert space $\cH$, $\conv H$ is the convex hull of $H$. 

One other particularly interesting set of states is the \emph{separable} states. We adopt the following notation for the set of density operators regarding separable states,
\begin{align*}
  \SEP(d,r) \coloneqq \conv{\psi_1 \otimes \cdots \otimes \psi_r \mid \ket{\psi_1},\ldots,\ket{\psi_r} \in \C^d}.
\end{align*}
A related notion is that of separable measurement, whose formal definition is given below.
\begin{definition}[Separable measurement]
A measurement $M = (M_0, M_1)$ is separable if in the yes case, the corresponding positive semi-definite matrix $M_1$ can be represented as a conical combination of two operators acting on the first and second parts, i.e.,
%has the following representation 
for some distribution $\mu$ over the tensor product of positive semi-definite matrices $\alpha$ and $\beta$ on the corresponding space,
\[
M_1 = \int  \alpha \otimes \beta ~\mathrm{d}\mu.
\]
\end{definition}
We record the following well-known fact. An interested reader is referred to~\cite{harrow2013church} for a formal proof.
\begin{fact}[Folklore] \label{fact:swap-test-sep}
The swap test is separable.
\end{fact}



\paragraph{Matrix Analysis} Given any matrix $M\in \C^{n\times n}$, $M^\dagger$ is its  conjugate transpose. Let $\sigma_1 \ge \sigma_2 \ge \ldots \ge \sigma_n$ denote its singular values. Then the trace norm $\|\cdot\|_1$, Frobenius norm $\|\cdot\|_F$ are defined as below
\begin{align*}
    \|M\|_1 = \sum_i \sigma_i, \qquad \|M\|_F = \sqrt{\sum_i \sigma_i^2} .
\end{align*}
The Frobenius norm also equals the square root of sum of squared modulus of each entry, i.e., $\|M\|_F=\sqrt{\sum_{i,j} |M(i,j)|^2}$.

For a positive semi-definite (PSD) matrix $M$, $\|M\|_F = \sqrt{\trace M ^2}$.  
For two PSD matrices, there is one (of many) analogous matrix Cauchy-Schwarz inequality.
\begin{equation}\label{eq:matrix-cauchy-schwarz}
\trace(\sigma\rho)\le \|\sigma\|_F \cdot\|\rho\|_F.
\end{equation}
We adopt the notation $\succeq$ to denote the partial order  that $\sigma \succeq \rho$ if $\sigma-\rho$ is  positive semi-definite.

\paragraph{Distances between Quantum States}
A standard notion of distance for quantum states is that of the \emph{trace distance}. The trace distance between $\psi$ and $\phi$, denoted $\TD(\psi, \phi)$, is 
\begin{align}
\frac{1}{2}\|\psi-\phi\|_1 = \frac{1}{2}\trace\sqrt{( \psi - \phi)^\dagger (\psi-\phi)}.\label{eq:trace-dist}
\end{align}
We also use the notation $\TD(\ket\psi, \ket\phi)$ if we want to emphasize that $\psi$ and $\phi$ are pure states. The following fact provides an alternative definition for trace distance between pure states.
\begin{fact}\label{fact:trace-innerproduct}
  The trace distance between $\ket{\phi}$ and $\ket{\psi}$ is given by $\TD(\ket{\phi},\ket{\psi}) = \sqrt{1-\abs{\braket{\phi}{\psi}}^2}$.
\end{fact}

Two states with small trace distance are indistinguishable to quantum protocols.
\begin{fact}\label{fact:trace_norm_acc}
If a quantum protocol accepts a state $\phi$ with probability at most $p$, then it accepts $\psi$  with
  probability at most $p + \TD(\phi,\psi)$.
\end{fact}


Trace distance enjoys the triangle inequality. For pure states, we can actually strengthen it.

\begin{claim} \label{fact:overlap-triangle-ineq}
Given unit vectors $\ket\alpha, \ket\phi, \ket\beta \in \cH$ for some Hilbert space $\cH$. Suppose
    \begin{align*}
        |\langle \alpha \mid \phi \rangle|^2 = 1-\epsilon, \qquad
        |\langle \beta \mid \phi \rangle|^2 = 1-\delta. 
    \end{align*}
Then for any $\epsilon + \delta \le 1$,\footnotemark
    \begin{equation}      
    |\langle \alpha \mid \beta \rangle|^2 \ge 
    (
        \sqrt{(1-\epsilon)(1-\delta)}-\sqrt{\epsilon\delta}
    )^2.\label{eq:overlap-triangle-ineq-tight}
    \end{equation}
In general, we always have 
    \begin{equation}      
    |\langle \alpha \mid \beta \rangle|^2 \ge 1-\epsilon-\delta-2\sqrt{\epsilon\delta}. \label{eq:overlap-triangle-ineq-crude}
    \end{equation}
\end{claim}
\footnotetext{When $\epsilon+\delta > 1$, then $\ket\alpha$ and $\ket\beta$ in general can be orthogonal.}
\begin{proof}
Without loss of generality assume that
\begin{align}
    &\ket\alpha = \sqrt{1-\epsilon}\ket\phi + \sqrt\epsilon \ket \mu, 
    \nonumber\\
    &\ket\beta = \sqrt{1-\delta}\ket\phi + \sigma\sqrt\eta\ket\mu + \sqrt{\delta-\eta}\ket\rho,
    \nonumber
\end{align}
where $\ket\mu, \ket\rho, \ket\phi$ are orthogonal, $0\le\eta\le \delta$ and $\sigma\in \C$ is a relative phase. Using the basis $\{\ket\phi, \ket\mu, \ket\rho\}$, we can write down explicitly the density matrix of $\alpha$ and $\beta$:
\begin{align*}
&\alpha = \begin{pmatrix}
    1-\epsilon & \sqrt{\epsilon(1-\epsilon)} & 0 \\
    \sqrt{\epsilon(1-\epsilon)} & \epsilon & 0\\
    0 & 0&0
\end{pmatrix},
%
\\
&\beta = \begin{pmatrix}
    1-\delta   &   \sigma\sqrt{(1-\delta)\eta}  &  
 \sqrt{(1-\delta)(\delta-\eta)} \\
    \sigma^* \sqrt{(1-\delta)\eta} &   \eta    &   \sigma\sqrt{\eta(\delta-\eta)}\\
    \sqrt{(1-\delta)(\delta-\eta)} &   \sigma^*\sqrt{\eta(\delta-\eta)} & \delta-\eta
\end{pmatrix}.
\end{align*}
Now by definition, 
\begin{align}
    \TD(\ket\alpha, \ket\beta) ^2
    &= \left(\frac{1}{2}\trace\sqrt{(\alpha - \beta)^\dagger (\alpha-\beta)}\right)^2
       \nonumber \\
    &= \frac{1}{2} \|\alpha-\beta\|_F^2
        \nonumber\\
    &= \frac{1}{2}( (\epsilon-\delta)^2 + (\epsilon-\eta)^2 + (\delta-\eta)^2 )
        +\eta(\delta-\eta) 
        \nonumber \\
    &\qquad\qquad + |\sqrt{\epsilon(1-\epsilon)}-\sigma\sqrt{(1-\delta)\eta}|^2 +(1-\delta)(\delta-\eta) 
        \nonumber\\
    &\le \frac{1}{2}( (\epsilon-\delta)^2 + (\epsilon-\eta)^2 + (\delta-\eta)^2 )
        +\eta(\delta-\eta) 
        \nonumber \\
    &\qquad\qquad + (\sqrt{\epsilon(1-\epsilon)}+\sqrt{(1-\delta)\eta})^2 +(1-\delta)(\delta-\eta), \label{eq:trace-triangle-bound}
\end{align}
where the second step holds because $\alpha-\beta$ is Hermitian with trace 0 and rank 0 or 2.
We claim that the RHS of (\ref{eq:trace-triangle-bound}), denote by $f$, is non-decreasing for $\eta\in[0,\delta]$. 
%We verify that $d/d\eta \,\TD(\ket\alpha, \ket\beta)$ is non-negative for $0\le\eta\le \delta$:
By routine calculation, 
\begin{align*}
    \frac{d f}{d\eta}  = -\epsilon + \sqrt{ \frac{\epsilon}{\eta} (1-\epsilon)(1-\delta)} \ge 0 \quad
    \iff \quad(1-\epsilon)(1-\delta) \ge \eta\epsilon \quad
    \\
    %
    \impliedby 
    \quad
    (1-\epsilon)(1-\delta) \ge \delta\epsilon
    \quad
    %
    \iff \quad 1 \ge \epsilon + \delta.
\end{align*}
As we assumed that $1 \ge \epsilon+\delta$, $df/d\eta$ is always non-negative.
Since the RHS of (\ref{eq:trace-triangle-bound}) is non-decreasing for $\eta\in[0,\delta]$, plug $\eta=\delta$ into the RHS of (\ref{eq:trace-triangle-bound}), we obtain
\begin{align*}
    \TD(\ket\alpha,\ket\beta) ^2
    &\le (\epsilon-\delta)^2 + (\sqrt{\epsilon(1-\epsilon)}+\sqrt{(1-\delta)\delta})^2,
\end{align*}
In view of \cref{fact:trace-innerproduct}, (\ref{eq:overlap-triangle-ineq-tight}) is proved. The ``in general'' part is trivially true when $\epsilon+\delta>1$ and otherwise follows from (\ref{eq:overlap-triangle-ineq-tight}).
\end{proof}

Another widely used distance measure between quantum states is that of \emph{fidelity}. For any density operators $\rho, \sigma$ from the same Hilbert space,
\begin{equation*}
    F(\rho, \sigma) = \left(\trace \sqrt{\sqrt\rho \sigma \sqrt\rho)}\right)^2.
\end{equation*}
For our purposes, we only need the fact that when one of the two density operators corresponds to a pure state, then
\begin{equation*}
    F(\rho,\sigma) = \trace(\rho\sigma).
\end{equation*}
The well-known data processing inequality for fidelity states that applying quantum operation never decreases the fidelity. 
\begin{fact}\label{fact:data-processing}
For any quantum channel (CPTP map) $\Phi$,
    \[F(\Phi(\rho), \Phi(\sigma)) \ge F(\rho, \sigma).\]
\end{fact}





\paragraph{Schmidt Decomposition and Partial Trace} 
For $\ket\psi$ describing quantum states over two subsystems $A,B$, e.g., $\ket\psi\in\C^m \otimes \C^n$, there are two sets of orthonormal states $\{\ket{\alpha_1}, \ket{\alpha_2}, \ldots, \ket{\alpha_k}\}\subseteq \C^m$, $\{\ket{\beta_1}, \ket{\beta_2}, \ldots, \ket{\beta_k}\}\subseteq \C^n$ , and positive numbers $\lambda_1\ge \lambda_2\ge \cdots \ge \lambda_k$ for some $k\le \min\{n,m\}$ such that
\begin{equation}
    \ket\psi = \sum_{i=1}^k \sqrt{\lambda_i} \ket{\alpha_i}\ket{\beta_i}, \qquad \text{and } \sum_{i=1}^k \lambda_i = 1.\label{eq:schmidt-decomp}
\end{equation}
The formula (\ref{eq:schmidt-decomp}) is called the \emph{Schmidt decomposition} of $\ket\psi$. The set of $\sqrt{\lambda_i}$ is unique, and is called the \emph{Schmidt coefficient} of $\ket\psi$. 
We  call $\sqrt{\lambda_1}$ the \emph{top Schmidt coefficient} and $\ket{\alpha_1}\ket{\beta_1}$ the \emph{top Schmidt component}. 
Note that the top Schmidt component may not be unique ignoring the global phases, in that case we break  tie arbitrarily. Since Schmidt decomposition follows from singular value decomposition, the (top) Schmidt coefficients can also be formulated as some optimization problem.
\begin{claim}\label{claim:top-schmidt-coeff}
Given any state $\ket\psi\in \cH_1\otimes\cH_2$. Then
\begin{equation*}
    \lambda_1 = \max_{\ket\sigma\in \cH_1, \ket\rho\in\cH_2}|\langle \psi \mid \sigma,\rho \rangle |^2
\end{equation*}
\end{claim}

Often we want to study the density operator of a quantum state $\ket\psi$ over the subsystem $A$, mathematically described by tracing out $B$, denoted  $\trace_B(\psi)$.
We also abbreviate $\psi_A=\trace_B(\psi)$. Note that fidelity never increases under partial trace due to \cref{fact:data-processing}, and similarly, the trace distance never increases under partial trace:
\begin{fact} \label{fact:trace-dist-partial-trace}
For any quantum states $\psi$ and $\phi$ over systems $A$ and $B$,
\[
    \TD(\psi, \phi) \ge \TD(\psi_A, \phi_A).
\]
\end{fact}
We use subscript to emphasize the systems that an operator is describing, e.g., $\psi^{AB}$ simply means that $\psi$ is a state over systems $A$ and $B$.

\iffalse
For a logical condition $C,$ we use the Iverson bracket
\[
\1{C}=\begin{cases}
1 & \text{if \ensuremath{C} holds,}\\
0 & \text{otherwise.}
\end{cases}
\]
\fi


\paragraph{Quantum Merlin-Arthur Systems}
We now formally define the class $\textup{almost-QMA}^{\R}(k)$, but
first we will need the $\ell_2$-sign bias definition, which, roughly
speaking, quantifies the imbalance in $\ell_2$ mass between the
positive and negative amplitudes parts of a state.

\begin{definition}[$\ell_2$-sign bias]\label{def:ell2_sign_bias}
  Given $\ket{\psi} \in \R^n$, we can uniquely write it as $\ket{\psi}= \sqrt{a} \ket{\psi_+} + \sqrt{1-a} \ket{\psi_{-}}$,
  where $a \in [0,1]$, $\ket{\psi_{+}}$ and $\ket{\psi_{-}}$ are unit vectors with only positive and negative amplitudes, respectively.
  The $\ell_2$-sign bias of $\ket{\psi}$ is defined as $\abs{a-(1-a)}$.
\end{definition}

  Note that a non-negative amplitude state has $\ell_2$-sign bias of $1$ whereas a general
  state has bias at least $0$.
  \textup{Almost}-$\QMA^\R(k)$ will be defined based on $\ell_2$-sign as a natural relaxation of $\QMA^+(k)$ towards the general $\QMA(k)$. 


\begin{definition}[$\textup{almost-QMA}^\R(k)$]\label{def:almost_qma_k}
  Let $k \colon \mathbb{N} \to \mathbb{N}$ be a polynomial time computable function.
  %
  A promise problem $\mathcal{L}_{\textup{yes}},\mathcal{L}_{\textup{no}} \subseteq \set{0,1}^*$ is in $\textup{almost-QMA}^\R(k)$ if there exists a $\BQP$
  verifier $V$ such that for every $n \in \mathbb{N}$ and every $x \in \set{0,1}^n$, 
  \begin{itemize}[label={}]
  \item \textbf{\textup{Completeness:}} If $x \in \mathcal{L}_{\textup{yes}}$, then there exist unentangled states $\ket{\psi_1}, \ldots, \ket{\psi_{k(n)}}$, each of $\ell_2$-sign
    bias $1/\poly(n)$ and on at most $\poly(n)$ qubits,
    s.t.\ $\Pr[V(x, \ket{\psi_1} \otimes \cdots \otimes \ket{\psi_{k(n)}})\textup{ accepts}] \ge 9/10$.
\vspace{\parsep}
    \item \textbf{\textup{Soundness:}} If $x \in \mathcal{L}_{\textup{no}}$, then for every unentangled states $\ket{\psi_1}, \ldots, \ket{\psi_{k(n)}}$, each of each of $\ell_2$-sign
    bias $1/\poly(n)$ and on at most $\poly(n)$ qubits,
    we have \ $\Pr[V(x, \ket{\psi_1} \otimes \cdots \otimes \ket{\psi_{k(n)}})\textup{ accepts}] \le 1/10$.
  \end{itemize}
\end{definition}

